\begin{code}[hide]
module paper.sections.examples where

open import Algebra.Bundles using (Group ; Monoid)
open import Algebra.Morphism.Structures using (module MonoidMorphisms)
open import Data.Nat using (ℕ ; suc ; 2+)
open import Data.Nat.Primality using (Prime)
open import Level using (0ℓ)

open import Notations using (₁₊)
open import Word.Base using (_ʷ)
open import Presentation.Definitions using (_IsPresentationOf_)
open import Presentation.Construct.Base using (_⋄_⋄_ ; CommRel ; ConjRelʷ ; _⊕^_)
import Presentation.Properties as PP
import Presentation.Construct.Properties.DirectProduct as DP
import Presentation.Construct.Properties.NDirectProduct as NDP
open import Examples.Groups.Cyclic.Normalization using (_Cn,_===_)
open import Examples.Groups.Cyclic.Semantics using (Cn-group)
import Examples.Groups.Cyclic.Presentation as CyP
open import Examples.Groups.Symmetric.Syntactics using (_VRel,_===_)
import Examples.Groups.Pauli.Presentation as Pauli
import Examples.Construct.SemiDirectProduct.SnD as SnD
import Examples.Amalgamations.CliffordT1 as CliffordT1

open SnD using (conj ; module Wreath)
open Pauli using (Pauli-group) renaming (Γ-H to Γ-Pauli₁)
open Monoid using (rawMonoid)
open MonoidMorphisms using (IsMonoidIsomorphism)
open PP using (•-ε-monoid)
open CliffordT1.CliffordT1 using (_===_ ; amalPres ; f)
\end{code}
\section{Examples: A Catalogue of Verified Presentations}\label{sec:examples}

This section is the library's harvest.  Table~\ref{tab:presentations} lists
every concrete presentation theorem in the development
(\aid{{\_}IsPresentationOf{\_}}, or its monoid version for the order-$0$
cyclic case);
Table~\ref{tab:nfs} lists the normal-form and uniqueness witnesses behind
them.  All of them are restated with full types in the index module
\texttt{MainTheorems.agda}, which also re-exports the conditional
construction theorems of \S\ref{sec:constructions} with their premise
telescopes.  We then discuss each family, ending with the three
amalgamation case studies, whose final theorems are monoid
\emph{isomorphisms} between gate-set presentations and amalgamated
products of presentations.  Throughout, $\bot$ is the empty type and $\top$ the unit
type, whose sole element is \aid{tt}; used as alphabets they have no
generator and a single generator respectively, and we write $T$ for the
one-letter word $[\,\aid{tt}\,]$\aid{ʷ}.

\begin{table}
\caption{Concrete presentation theorems (all verified; names as in
\texttt{MainTheorems.agda}).  $p$ ranges over primes, $n, m, N$ over
naturals.}
\label{tab:presentations}
\centering
\small
\begin{tabular}{@{}lll@{}}
\toprule
Relation set & Presents & Index name \\
\midrule
\aid{EmptyRel} over $\bot$ & trivial group & \aid{trivial-presentation} \\
\aid{TrivialRel} over any $A$ & trivial group & \aid{trivial-presentation′} \\
$T^{\,n+1} = \varepsilon$ & $\mathbb{Z}/(n{+}1)\mathbb{Z}$ & \aid{cyclic-presentation} \\
$T^{\,0} = \varepsilon$ & the \emph{monoid} $(\mathbb{N}, +, 0)$ & \aid{cyclic-monoid-presentation} \\
\aid{order}, \aid{yang-baxter} + structural & permutations of $\mathsf{Fin}\;n$ & \aid{symmetric-presentation} \\
$C_p \aid{⋄} C_p \aid{⋄} \aid{CommRel}$ & $\mathbb{Z}_p \times \mathbb{Z}_p$ & (\aid{H-pres}, Pauli development) \\
$(C_p \aid{⋄} C_p \aid{⋄} \aid{CommRel})\,\aid{⊕\char94}\,n$ & Pauli quotient $(\mathbb{Z}_p{\times}\mathbb{Z}_p)^n$ & \aid{pauli-presentation} \\
$(C_N\,\aid{⊕\char94}\,n) \aid{⋄} S_n\text{-circuits} \aid{⋄} \aid{ConjRelʷ}$ & $\mathbb{Z}_N \wr S_n$ (generalized symmetric groups) & \aid{wreath-presentation} \\
\bottomrule
\end{tabular}
\end{table}

\begin{table}
\caption{Normal-form and uniqueness witnesses (selection).  NFI =
\aid{NormalFormInjective}, NF = \aid{NormalForm}, UNF =
\aid{UniqueNormalForm}.  For the constructions, UNF is lifted from the
factor witnesses.}
\label{tab:nfs}
\centering
\footnotesize
\setlength{\tabcolsep}{4pt}
\begin{tabular}{@{}llll@{}}
\toprule
Presentation & Carrier & Witnesses & Note \\
\midrule
cyclic, order $n{+}1$ & $\mathsf{Fin}\,(n{+}1)$ & NFI, NF, UNF & presents $\mathbb{Z}/(n{+}1)\mathbb{Z}$ \\
cyclic, order $0$ & $\mathbb{N}$ & NFI, NF, UNF & presents the monoid $(\mathbb{N}, +, 0)$ (\S\ref{sec:cyclic}) \\
$S_n$ circuits & $\top \times C\,1 \times \dots \times C\,(n{-}1)$ & NFI, NF, UNF$\times 2$ & factorial system, $|{\cdot}| = n!$ \\
$\Gamma \aid{⋄} \Delta \aid{⋄} \aid{CommRel}$ & $\mathit{NF}_1 \times \mathit{NF}_2$ & NFI, NF, UNF & lifted from factors \\
$\Gamma \aid{⋄} \Delta \aid{⋄} \aid{ConjRelʷ}\,\mathit{conj}$ & $\mathit{NF}_1 \times \mathit{NF}_2$ & NFI, NF, UNF & twist on the left factor \\
$P_1 \mathrel{\aid{*}} P_2 \mathrel{\aid{⋆}} f_1 \mathrel{\aid{⋆}} f_2$ & $\mathit{NF}_0 \times \aid{amalNFC}\;C\;D$ & NFI, NF, UNF & base NF, then alternating cosets (\S\ref{sec:constructions}) \\
Clifford+$T$ base $\langle \omega \rangle {\times} \langle S \rangle$ & $\mathsf{Fin}\,8 \times \mathsf{Fin}\,4$ & NFI, NF, UNF & radices $8, 4$; unique for $\mathbb{Z}_8 {\times} \mathbb{Z}_4$ \\
\bottomrule
\end{tabular}
\end{table}

\subsection{Cyclic groups, and the trivial group twice}\label{sec:cyclic}

The cyclic development presents $\langle\, \top \mid T^{\,n+1} =
\varepsilon\,\rangle$ with normal-form carrier $\mathsf{Fin}\,(n{+}1)$,
for every positive order at once (\aid{cyclic-presentation}): the normal
form of a word is its letter count modulo $n{+}1$, the section sends the
residue $k$ to $T^{k}$, and uniqueness for
$\mathbb{Z}/(n{+}1)\mathbb{Z}$ is immediate.  Order $0$ gets a theorem
of its own, \aid{cyclic-monoid-presentation}: with carrier $\mathbb{N}$
(a word count with no wraparound), $T^{\,0} = \varepsilon$
\aid{IsMonoidPresentationOf} $(\mathbb{N}, +, 0)$, the free monoid
rather than the group $\mathbb{Z}$, since, as \S\ref{sec:wordproblem}
explained, the generator has no inverse there and the \aid{Grouplike}
upgrade is exactly what fails.  The
trivial group is presented twice (empty alphabet with no relations, and
arbitrary alphabet with the everything-is-$\varepsilon$ relation), and
the two presentations are proved isomorphic.

\subsection{Wreath products: signed permutations and beyond}\label{sec:wreath}

The headline compositional result combines almost every construction in
the library.  Fix a base order $N = m{+}1$ and $n$ wires.  The $n$-fold
direct power $(C_N)^{\oplus n}$ presents $(\mathbb{Z}_N)^n$; the circuit
presentation of $S_n$ acts on it by permuting coordinates, an action
specified \emph{syntactically} as a word-valued conjugation
$\mathit{conj} : \aid{Gen}\;n \to (\top \uplus^n) \to
\aid{Word}\,(\top \uplus^n)$.  Two finite compatibility lemmas (the action
respects the Coxeter relations pointwise, and respects the product
relations) feed the semidirect-product presentation theorem, giving
\begin{code}
wreath-presentation : ∀ n m → (((suc m Cn,_===_) ⊕^ n) ⋄ (n VRel,_===_) ⋄ ConjRelʷ conj)
                                IsPresentationOf (Wreath.wreath-group n m)
\end{code}
\begin{code}[hide]
wreath-presentation n m = Wreath.presentation n m
\end{code}
where $\aid{wreath-group}\;n\;m$ is the semantic
$(\mathbb{Z}_N)^n \rtimes S_n = \mathbb{Z}_N \wr S_n$ built from the
\texttt{ForStdlib} semidirect product.
The lifted normal form is a pair of $n$ residues modulo $N$ and the
factorial tuple of a permutation, so its section is a word of $T$-powers,
one per coordinate, followed by the staircases of
\S\ref{sec:permutations}.  At $N = 2$ this is the group of
\emph{signed permutations}, the hyperoctahedral group $B_n$; at
general $N$ it is the generalized symmetric group $G(N,1,n)$ of monomial
matrices with $N$-th-root-of-unity entries.  These groups are ubiquitous
in the qudit setting (Pauli-frame changes, permutation-and-phase
subgroups of Clifford groups), and the theorem delivers their
presentations \emph{for all $N$ and $n$ at once}: a family of
completeness results parameterized over two naturals, obtained
compositionally rather than by a separate argument per instance.

\subsection{Pauli groups modulo phases in all prime dimensions}\label{sec:pauli}

The Pauli development is the purest showcase of compositionality: there
is \emph{no} coset table in it at all:

\begin{code}
pauli-presentation : ∀ m (prime : Prime (2+ m)) n →
  (Γ-Pauli₁ m prime ⊕^ n) IsPresentationOf (Pauli-group m prime n)
\end{code}
\begin{code}[hide]
pauli-presentation m prime = Pauli.Pauli-presentation m prime
\end{code}

For a prime $p = 2 + m$, with \aid{prime} its primality proof (a
\emph{qupit} is a qudit whose dimension is a prime), the single-qupit
relation set $\aid{Γ-Pauli₁}$ (\aid{Γ-H} in the index module) is the
direct-product
join of two order-$p$ cyclic presentations ($X$- and $Z$-type generators,
which commute in the quotient of the Pauli group by its phase subgroup),
and the $n$-qupit relation set is its $n$-fold power.  The semantic group
$(\mathbb{Z}_p \times \mathbb{Z}_p)^n$ is assembled by the standard
library's binary direct product and our $n$-fold iteration, and the
presentation is literally three construction theorems composed (the
cyclic presentation, the binary product, and the $n$-fold product).  The
normal form is a tuple of $n$ pairs of exponents in
$\mathsf{Fin}\,p \times \mathsf{Fin}\,p$, with section
$X^{a_1} Z^{b_1} \cdots X^{a_n} Z^{b_n}$.  The
full Pauli group with phases, and its role beneath the Clifford
hierarchy, is left to future work.

\subsection{Single-qubit Clifford+$T$, as an amalgamated free product}\label{sec:cliffordt1}

The first amalgamation case study presents the single-qubit Clifford+$T$
\emph{monoid}, the flagship example of the completeness literature for
good reasons: Clifford+$T$ is the smallest universal gate set in common
use, its single-qubit fragment has the first and best-known canonical
form (the Matsumoto--Amano normal
form~\cite{matsumoto2008representation,giles2013remarks}, which underlies
exact synthesis and $T$-count optimality), and the later normal forms
for single-qutrit and single-qudit
Clifford+$T$~\cite{glaudell2019qutrit,prakash2018qutrit,jain2020qudit}
follow its syllable pattern.  The relation set (over generators
$T, H, S, \omega$, with $X := H S S H$ as an abbreviation) is
\[
  \omega^8 = \varepsilon,\quad
  T^2 = S,\quad
  S^4 = \varepsilon,\quad
  H^2 = \varepsilon,\quad
  (SH)^3 = \omega,\quad
  (TX)^2 = \omega,\quad
  \omega\ \text{central},
\]
and the verified structure exhibits the presented monoid, written
$\mathrm{Clifford}{+}T$ below (its identification with the operator
monoid is the unformalized step discussed in \S\ref{sec:threats}), as an
amalgamated free product over the subgroup $M$ generated by $X$, $S$, and
$\omega$:
\[
  \mathrm{Clifford}{+}T \;\cong\;
  \langle T, M \rangle \,\ast_{M}\, \langle H, M \rangle,
  \qquad M = \langle X, S, \omega \rangle .
\]
The development builds a tower whose base is the direct product
$\langle \omega \rangle \times \langle S \rangle$ of cyclic groups of
orders $8$ and $4$ (normal-form radices $8 \cdot 4$), extends it by one
verified coset table to the subgroup $M$ (two cosets:
$\varepsilon$ and $X$), and then extends \emph{twice from $M$}:
by a two-coset table (cosets $\varepsilon, T$) to
$\mathit{TXS\omega} = \langle T, X, S, \omega \rangle$, and by a
three-coset table (cosets $\varepsilon, H, HS$) to the Clifford group
$\langle H, X, S, \omega \rangle$.  The two \aid{PackedCosetTable}s over
the shared $M$ form an \aid{AmalDataNF}, and the
amalgamation module produces the normal form for the amalgamated
product.  It is instructive to see what this carrier
\emph{is}: with a one-element $C$ ($T$-coset) and two-element $D$
($H, HS$), a normal form is an element of $\mathit{NF}_0 \times
\aid{amalNFC}\;C\;D$, and its section spells it out, in the library's
left-to-right word order, as the regular expression
\[
  \mathcal{M}\;(\varepsilon \mid H \mid HS)\;(TH \mid THS)^{*}\;(\varepsilon \mid T),
\]
where $\mathcal{M}$ stands for the base normal form of the $M$-part; an
element of $\aid{amalNFC}\;C\;D$ is an optional leading
$H$-syllable, a list of alternations, and an optional trailing $T$,
which, read right to left, is the Matsumoto--Amano normal
form~\cite{matsumoto2008representation,giles2013remarks},
with its single Clifford factor $\mathcal{C}$ split into a coset
representative and an element of $M$:
the alternating normal form of an amalgamated free product, as observed
for this group by Bian and Selinger~\cite{bian2023cliffordcs3}.  The
final theorem is the
syntactic isomorphism assembled by \aid{StarIsomorphism} from
generator-level translations both ways (with \aid{{\_}==={\_}} the
gate-set relations, \aid{amalPres} the amalgamated presentation, and
\aid{f} the translation):

\begin{code}
clifford+T-qubit-isomorphism : IsMonoidIsomorphism
  (rawMonoid (•-ε-monoid _===_)) (rawMonoid (•-ε-monoid amalPres)) (f ʷ)
\end{code}
\begin{code}[hide]
clifford+T-qubit-isomorphism = CliffordT1.CliffordT1.CliffordT1-isomorphism
\end{code}

Every well-definedness obligation in the tower (on all axioms at all
cosets, 59 cases) is discharged by
\aid{(by-equal-nf Eq.refl) , Eq.refl}: the typechecker runs the table
on both sides of the axiom from the given coset, the resulting cosets
coincide, and the residual words have the same normal form one level
down.

\subsection{Single-qutrit Clifford+$T$}\label{sec:qutrit}

The qutrit analogue replaces $\omega$ by a ninth root $\zeta$ and runs a
longer tower: from the cyclic base $\langle \zeta \rangle$ through
$\langle S, \zeta \rangle$ (three cosets), then a
\emph{nine}-coset table to $\langle S, X, \zeta \rangle$
(with the twisted commutation
$(XS)(SX) = \zeta^6 (SX)(XS)$), then a two-coset table adjoining the
squared Hadamard $H\!H$, and finally an amalgamation of a $T$-side and an
$H$-side extension over the common subgroup.  The presented gate-set
relations (generators $T, H, S, \zeta$; abbreviations
$X := \zeta^3 (HSH)(HSSH)$ and $Z := \zeta^3 S^2 H^2 S H^2$):
\[
  \zeta^9 = \varepsilon,\;\;
  S^3 = \zeta^6,\;\;
  H^4 = \varepsilon,\;\;
  T^3 = Z,\;\;
  (SH)^3 = \varepsilon,\;\;
  (THH)^2 = \varepsilon,
\]
\[
  HHSHHS = SHHSHH,\;\;
  TS = ST,\;\;
  TX = \zeta^3 S^2 X\, T,\;\;
  \zeta\ \text{central},
\]
a presentation derived and verified in this development.  The verified
structure is again an amalgamated free product of the presented monoid,
\[
  \mathrm{Clifford}{+}T_3 \;\cong\;
  \langle T, M \rangle \,\ast_{M}\, \langle H, M \rangle,
  \qquad M = \langle S, X, \zeta, H^2 \rangle.
\]
The alternating normal form reads
\[
  \mathcal{M}\;(\varepsilon \mid H \mid HS \mid HS^2)\;
  \bigl((T \mid TH^2)(H \mid HS \mid HS^2)\bigr)^{*}\;
  (\varepsilon \mid T \mid TH^2),
\]
and right to left it is the $T$-optimal canonical form of Glaudell,
Ross, and Taylor,
$(\varepsilon \mid T \mid H^2T)\,(HT \mid H^3T \mid SHT \mid SH^3T \mid
S^2HT \mid S^2H^3T)^{*}\,\mathcal{C}$~\cite{glaudell2019qutrit}
(Prakash, Jain, Kapur, and Seth independently found a similar
form~\cite{prakash2018qutrit}, extended to qudits
in~\cite{jain2020qudit}), with $\mathcal{C}$ split as in the qubit
case.  The final theorem has the same shape as
the qubit one (\aid{clifford+T-qutrit-isomorphism} in the index)
and the same proof economy, the nine-coset level's seventy-two
axiom-preservation obligations all closing by
\aid{refl} after evaluation (\S\ref{sec:engine}).  We know of no prior
machine-checked normal form for a qutrit gate-set presentation.

\subsection{\texorpdfstring{$U_3(\mathbb{Z}[\tfrac12,i])$}{U3(Z[1/2,i])}}\label{sec:u3}

The third case study formalizes the case $n = 3$ of the presentation, due
to Bian and Selinger, of the groups
$U_n(\mathbb{Z}[\tfrac12,i])$~\cite{bian2021un}, the unitary matrices over
the ring $\mathbb{Z}[\tfrac12,i]$.
Generators are $i_0$ (a local phase), $K_{01}$ (a scaled $2{\times}2$
Hadamard-type block), and the two transpositions $X_{01}, X_{12}$;
derived generators $i_1, i_2, K_{02}, K_{12}$ are conjugates.  The eleven
relations (labelled $\mathrm{S1}$--$\mathrm{S14}$ in the development, a
numbering with gaps) include the orders $i_0^4 = X_{01}^2 = X_{12}^2 =
\varepsilon$, braid and commutation laws for the transpositions and
phases, and the interactions of $K_{01}$ with the
phases, e.g.
$K_{01} i_1^2 = X_{01} K_{01}$ and
$K_{01}^2\, i_0\, i_1 = \varepsilon$.  The verified structure, with
$U_3(\mathbb{Z}[\tfrac12,i])$ standing for the presented group
(\S\ref{sec:threats}), is
\[
  U_3(\mathbb{Z}[\tfrac12,i]) \;\cong\;
  (\mathbb{Z}_4 \wr S_3) \,\ast_{M}\, \langle M, K_{01} \rangle,
  \qquad M = (\mathbb{Z}_4 \wr S_2) \times \mathbb{Z}_4 .
\]
The right factor adjoins
$K_{01}$ to $M$: a six-coset Reidemeister--Schreier level presents
the $\langle K_{01}, i_0, i_1 \rangle$ subgroup over the product
of two order-$4$ cyclic groups, a direct product then adjoins
$\langle i_2 \rangle$, and $X_{01}$ enters as a word in $K_{01}$, $i_0$,
$i_1$.  The alternating normal form reads
\[
  \mathcal{M}\;(\varepsilon \mid K_{01} \mid K_{01} i_0)\;
  \bigl((X_{12} \mid X_{12}X_{01})(K_{01} \mid K_{01} i_0)\bigr)^{*}\;
  (\varepsilon \mid X_{12} \mid X_{12}X_{01}),
\]
with $\mathcal{M}$ the normal form of $M$.
The final theorem,
\aid{U₃-presentation-isomorphism}, is again a monoid isomorphism between
the flat eleven-relation presentation and the amalgamated-product
presentation.  At
$2\,868$ and $1\,394$ source lines, the qutrit and $U_3$ studies are the
largest verified objects in the library, and both consist mostly of
finite tables and \aid{refl}-discharged obligations, the shape of
proof this framework was designed to make routine.
